% !TeX root = ../初等数论笔记.tex
% 由主文档通过 \input 载入；请编译主文档。
\setcounter{chapter}{9} % 保留原书章号。
\chapter{不定方程与平方和}
\label{ch:diophantine}

一个实数方程可能有一条连续曲线上的解，限制为整数后却可能没有解、只有有限多解，或形成特殊的无限序列。本章综合使用因式分解、同余、互素性与参数化。对于勾股数和平方和，我们不仅要给出漂亮的公式，还要说明为什么每个解都能被这些公式或判据覆盖。

\section{整数方程的解题框架}

整数方程的第一步是明确未知量范围。整除与因式分解常能把无限搜索化为有限个约数组合；同余可以给出无解障碍；不等式用于界定取值；参数化用于描述全部解。整数解必在每个模数下给出同余解，因此某个模数下无解足以否定整数解，但通过有限个模数的检验并不足以证明整数解存在。

一种实用的思考顺序是先检查奇偶和小模数，排除明显障碍；再寻找可分解的乘积或能提取的最大公因数；若仍有无限多个候选，就考虑参数化、递降或大小估计。这个顺序只是解题策略，不是保证所有方程都可解的算法。证明无解时找到一个矛盾即可，证明有解要构造对象，求全部解则必须同时控制所有可能。

\begin{example}
求 $xy=x+y+3$ 的全部正整数解。
\end{example}
\begin{solve}
改写为 $(x-1)(y-1)=4$。正整数解必须有 $x,y\ge2$，因此因子对只有 $(1,4),(2,2),(4,1)$，得到 $(x,y)=(2,5),(3,3),(5,2)$。因子枚举已覆盖所有可能，所以答案完整。
\end{solve}

\section{勾股数的参数化}

满足 $x^2+y^2=z^2$ 的正整数三元组称为\emph{勾股三元组（Pythagorean triple）}；若 $\gcd(x,y,z)=1$，称其本原（primitive）。本原解的三个数两两互素，否则某个公共素因子也会整除第三个数。

\begin{theorem}
\label{thm:pythagorean}
交换两条直角边后，每个本原勾股三元组唯一写为
\begin{equation}
 x=u^2-v^2,\qquad y=2uv,\qquad z=u^2+v^2,
 \label{eq:pythagorean-parameters}
\end{equation}
其中 $u>v>0$，$\gcd(u,v)=1$，且 $u,v$ 奇偶性不同；规定 $x$ 为奇边、$y$ 为偶边。
\end{theorem}
\begin{proof}
两边不能均偶，也不能均奇，否则平方和模 $4$ 余 $2$。故可令 $x,z$ 奇、$y$ 偶。由
\[
 \frac{z+x}{2}\frac{z-x}{2}=\left(\frac y2\right)^2
\]
知左侧两因子乘积为平方。它们互素：公因数整除其和 $z$ 与差 $x$。因此分别为 $u^2,v^2$，且 $u>v>0$，得到所示公式。互素性给出 $\gcd(u,v)=1$，$z$ 为奇数又给出二者奇偶性不同；$u^2=(z+x)/2,v^2=(z-x)/2$ 保证唯一性。

反之，公式直接满足勾股等式。若某个奇素数同时整除 $u^2-v^2$ 与 $u^2+v^2$，则它整除 $u,v$，矛盾；两者又都是奇数。因此所构造的三元组本原。
\end{proof}

一般正整数解由本原解同时乘以正整数得到。取 $(u,v)=(2,1),(3,2)$，分别得到 $(3,4,5)$ 与 $(5,12,13)$。

\begin{example}[参数条件怎样保证不漏解]
求所有斜边不超过 $20$ 的本原勾股三元组，不区分两直角边的次序。
\end{example}
\begin{solve}
由 $z=u^2+v^2\le20$ 且 $u>v\ge1$，有 $2\le u\le4$。$u=2$ 时只能取 $v=1$，得 $(3,4,5)$；$u=3$ 时 $v=1$ 与 $u$ 同奇，应排除，$v=2$ 得 $(5,12,13)$；$u=4$ 时 $v=1$ 得 $(15,8,17)$，$v=2$ 不互素，$v=3$ 又使斜边为 $25>20$。故只有这三组。若不要求本原，还要加入这些三元组的正整数倍，并重新检查斜边上界。
\end{solve}

参数 $u,v$ 的互素与奇偶条件不是装饰：例如 $(u,v)=(3,1)$ 给出 $(8,6,10)$，虽满足勾股等式，却不是本原解。证明参数化的反向时应检查“确实满足方程”和“确实满足本原要求”两件事，不能只代入平方等式就结束。

\section{无穷递降法}

\emph{无穷递降法（infinite descent）}的要点是：若一组解存在，就能构造同类且某个正整数参数严格更小的解；这与良序原理矛盾。

\begin{theorem}
\label{thm:fermat-four}
方程 $x^4+y^4=z^2$ 没有正整数解。
\end{theorem}
\begin{proof}
递降的目标是把任意正解变成右侧平方根更小的正解。之所以研究 $x^4+y^4=z^2$ 而不是直接研究 $x^4+y^4=z^4$，是因为两次勾股参数化会自然产生右侧为平方的同类方程，从而能够闭合递降过程。

若有解，取 $z$ 最小的一组。它满足 $\gcd(x,y)=1$：若 $d>1$ 同时整除 $x,y$，则 $d^4\mid z^2$，由唯一分解有 $d^2\mid z$，约去 $d^4$ 就得到右侧参数为 $z/d^2$ 的更小解。由奇偶性可令 $x$ 奇、$y$ 偶。本原勾股三元组 $(x^2,y^2,z)$ 给出
\[
 x^2=u^2-v^2,\qquad y^2=2uv,\qquad z=u^2+v^2,
\]
其中 $u,v$ 互素且奇偶不同。若 $u$ 偶而 $v$ 奇，则 $x^2\equiv3\pmod4$，不可能；所以 $u$ 奇、$v$ 偶。由 $(y/2)^2=u(v/2)$ 及两因子互素，得 $u=s^2,v=2t^2$，于是
\[
 x^2+(2t^2)^2=s^4.
\]
这是本原勾股三元组：若素数同时整除 $x,t$，由 $x^2=s^4-4t^4$ 就也整除 $s$，与 $\gcd(u,v)=1$ 矛盾。再次参数化得
\[
 x=r^2-w^2,\qquad t^2=rw,\qquad s^2=r^2+w^2,
\]
其中 $r,w$ 互素且为正整数。乘积为平方迫使 $r=a^2,w=b^2$，因此 $a^4+b^4=s^2$。而 $s^2=u<z$，特别有 $s<z$，构造出了更小的正整数解，矛盾。
\end{proof}

因此 $x^4+y^4=z^4$ 也无正整数解，只需把右侧视为 $(z^2)^2$。这一论证只处理指数 $4$，并未证明所有指数的情形。

回看递降证明，需要逐项核对新构造的数仍为整数、仍为正数、仍满足同一种方程，并且选定参数严格变小。仅把数值“变小一些”或把方程变成另一类方程，都不足以产生矛盾。这里新解是 $(a,b,s)$，缩小的是右侧平方根，从 $z$ 变为 $s$；它恰好符合最初选择最小 $z$ 的标准。

\section{两平方和定理}

\begin{lemma}
\label{lem:prime-two-squares}
奇素数 $p$ 能写成两个整数的平方和，当且仅当 $p\equiv1\pmod4$。
\end{lemma}
\begin{proof}
必要性来自平方模 $4$ 的余数：奇数平方和只能由一奇一偶产生，因而余 $1$。

若 $p\equiv1\pmod4$，由式\eqref{eq:quadratic-supplements}可取 $t$ 使 $t^2\equiv-1\pmod p$。令 $h=\lfloor\sqrt p\rfloor$，考虑 $(h+1)^2>p$ 个整数对 $(u,v)$，其中 $0\le u,v\le h$，将它们按 $u-tv\pmod p$ 分入 $p$ 个盒子。两个不同的对有相同余数，作差得到不全为零的 $x,y$，满足 $|x|,|y|\le h$ 且 $x\equiv ty\pmod p$。于是
\[
 0<x^2+y^2\le2h^2<2p,\qquad p\mid x^2+y^2.
\]
只能有 $x^2+y^2=p$。
\end{proof}

\begin{theorem}[两平方和定理（sum of two squares theorem）]
\label{thm:two-squares}
正整数 $n$ 可写成 $x^2+y^2$（$x,y\in\mathbb Z$），当且仅当每个满足 $p\equiv3\pmod4$ 的素因子 $p$ 在 $n$ 的分解中具有偶数指数。允许 $x$ 或 $y$ 为零。
\end{theorem}
\begin{proof}
若这样的 $p$ 整除 $x^2+y^2$，由推论\ref{cor:prime-three-mod-four}，$p$ 同时整除 $x,y$，所以可从和式提出 $p^2$。反复进行，$p$ 的指数只能为偶数。

反之，$2=1^2+1^2$，模 $4$ 余 $1$ 的素数由引理\ref{lem:prime-two-squares}表示为两平方和，模 $4$ 余 $3$ 的素数的偶次幂本身为平方。最后使用恒等式
\begin{equation}
 (a^2+b^2)(c^2+d^2)=(ac-bd)^2+(ad+bc)^2
 \label{eq:two-square-product}
\end{equation}
逐个相乘即可得到 $n$ 的表示。
\end{proof}

例如 $65=5\cdot13$ 满足判据，并有 $65=1^2+8^2=4^2+7^2$；$21=3\cdot7$ 不满足判据，因而无表示。判据解决可表示性，完整枚举全部表示则是进一步的问题。

\begin{example}[从判据走向构造]
判断 $325$ 是否为两平方和，并由较小数的表示构造两组表示。
\end{example}
\begin{solve}
$325=5^2\cdot13$，没有模 $4$ 余 $3$ 的素因子，故可以表示。先用 $25=3^2+4^2$、$13=2^2+3^2$，代入式\eqref{eq:two-square-product}得
\[
 325=(3\cdot2-4\cdot3)^2+(3\cdot3+4\cdot2)^2=6^2+17^2.
\]
把第二个表示改为 $13=2^2+(-3)^2$，得到 $325=18^2+1^2$。还可用 $25=5^2+0^2$ 得到 $325=10^2+15^2$。这些构造验证了多组表示的存在；若要求全部表示，还应加上枚举范围或其他完备性论证。
\end{solve}

\section{四平方和专题（选学）}

\begin{theorem}[拉格朗日四平方和定理；此处不证]
每个非负整数都能写成四个整数的平方和。
\end{theorem}

该结论称为\emph{四平方和定理（four-square theorem）}，其完整证明可见参考资料\cite{mit-four-squares}。与两平方和相比，四个变量消除了可表示性的限制。例如 $7=2^2+1^2+1^2+1^2$，但它不是两平方和；它也不是三平方和，因为平方模 $8$ 只可能余 $0,1,4$，三个这样的余数不可能相加得到 $7$。这里引用的是存在性结论，并未给出高效求表示的算法。

\section{习题}

求整数方程的全部解，需要把每次变换的正反方向都说明。勾股参数化依赖互素乘积为平方时各因子为平方；递降依赖正整数不能无限下降；两平方和判据则将整体可表示性化为素因子的类型与指数条件。

\begin{enumerate}
 \item 求斜边为 $25$ 的全部正整数勾股三元组，不区分两条直角边的次序。
 \item 判断 $45,75,130$ 哪些可表示为两个整数的平方和，并给出可表示者的一组表示。
 \item 证明不存在面积为奇数的整数边长直角三角形。
 \item 求 $xy=2x+3y$ 的全部正整数解，并说明枚举为何完整。
 \item\label{ex:dio-factors} 求 $x^2-y^2=45$ 的全部非负整数解。
 \item\label{ex:dio-hypotenuse} 证明本原勾股三元组的斜边的每个素因子都模 $4$ 余 $1$。
 \item\label{ex:dio-two-square-list} 判断 $50,77,98,325$ 是否为两平方和，并给出可表示者的一组表示。
 \item\label{ex:dio-area} 证明整数边长直角三角形的面积一定被 $6$ 整除。提示：先处理本原三元组，再乘回公因数。
\end{enumerate}

